\documentclass[9pt]{extarticle}
\usepackage[a4paper,margin=1.6cm]{geometry}
\usepackage[T1]{fontenc}
\usepackage[utf8]{inputenc}
\usepackage{times}
\usepackage{microtype}
\usepackage{multicol}
\usepackage{enumitem}
\usepackage{array}
\usepackage{tabularx}
\usepackage{booktabs}
\usepackage{xcolor}
\usepackage[hidelinks]{hyperref}
\usepackage{titlesec}
\usepackage{parskip}
\usepackage{amsmath}
\usepackage{amssymb}
\usepackage[round,authoryear]{natbib}

% ---------- Section styling ----------
\titleformat{\section}
  {\normalfont\bfseries\large}{\thesection.}{0.5em}{}
\titlespacing*{\section}{0pt}{8pt}{4pt}
\titleformat{\subsection}
  {\normalfont\bfseries\normalsize}{\thesubsection}{0.5em}{}
\titlespacing*{\subsection}{0pt}{6pt}{3pt}
\setlist{nosep,leftmargin=*,labelsep=0.5em}

\newcommand{\BASESIM}{XCIST/CatSim}

\pagestyle{empty}

\begin{document}

\begin{center}
  {\Large\bfseries MetalSynth-Pelvis: Conditioned synthesis of osteosynthesis implants\\[2pt]
  and local metal artifacts in pelvic CT scan volumes using Multi-Window Latent Diffusion}\\[3pt]
  {\normalsize Kiara Alexandra Balc\'azar Santa Cruz --- UTEC \quad $\cdot$ \quad Advisor: V\'ictor Flores Benites}
\end{center}
\vspace{2pt}\hrule\vspace{6pt}

% #####################################################################
\section{Problem Statement}
% #####################################################################

Automatic pelvic computed tomography (CT) analysis and volumetric segmentation depend on sufficiently annotated data. Metal artifacts can obscure anatomical boundaries near implants, motivating evaluation of segmentation robustness in this region. CTPelvic1K \citep{liu2021ctpelvic1k} includes a CLINIC-metal subset whose limited bone annotation constrains supervised evaluation. Its name does not establish that every case contains pelvic osteosynthesis, nor that other subsets are free of metal or foreign objects. The local audit will therefore determine implant types, eligible training anatomy and annotated evaluation cases through three-dimensional and multiplanar review. The magnitude of peri-implant segmentation degradation must be measured in this study rather than assumed from the dataset description.

While generative models based on diffusion processes excel at synthesizing visually plausible medical anomalies, they structurally fail when applied to metallic implants. State-of-the-art lesion synthesis models, such as DiffTumor \citep{chen2024tumorsynthesis}, DiffBoost \citep{zhang2025diffboost}, CLAIM \citep{ramzan2026claim}, and LGESynthNet \citep{jacob2026lgesynthnet}, are designed around bounded inpainting for deformable biological tissues. They implicitly assume non-rigidity and strictly restrict intensity alterations to the inside of the object's mask. This is physically incompatible with metal implants, which are rigidly geometric and produce severe attenuation profiles (beam hardening, photon starvation, and high-amplitude streaking) that propagate globally far beyond the implant's boundaries, destroying the signal of the surrounding soft tissue.

Physics-based metal artifact simulation provides a reference for appearance generation. This thesis adopts the hybrid simulation protocol of \citet{peters2025hybrid}, implemented with CatSim within XCIST \citep{deman2007catsim,wu2022xcist}, as its physical comparison protocol. Their phantom validation supports the evaluated simulation settings; it does not establish perfect physical coherence or validate pelvic osteosynthesis in three dimensions. Their training data use randomized metal placement, whereas their scoring benchmark uses expert placement. The unresolved placement problem is therefore scalable, anatomically constrained placement, rather than scalable metal insertion itself.

Even if automated placement were achieved, mere anatomical containment is insufficient; a valid simulation must reflect true surgical admissibility. Traditional deterministic surgical planning pipelines \citep{liu2025pipeline} attempt to optimize a single "perfect" trajectory. Yet, real-world clinical data reveals a different reality: \citet{zwingmann2009navigated} established that 31--60\% of real-world iliosacral screw placements result in some degree of malposition or cortical breach, based on the clinical scale of \citet{smith2006iliosacral}. Training a diagnostic network solely on "perfect" deterministic trajectories will inevitably cause it to fail when encountering the true variance of clinical errors. Therefore, a generative sampler must not optimize for perfection, but instead sample from the true statistical distribution of clinical malpositions, bounded by operative safe zones \citep{ramadanov2025safezone} and pelvic density representations \citep{arand2019pelvicring}.

\textbf{Recognized Gap.} Current attempts to bridge the metal-artifact domain shift are bottlenecked by three mutually exclusive paradigms: physical simulators do not supply scalable anatomically constrained placement, biological diffusion models cannot generate rigid global streaks, and deterministic surgical planners ignore the statistical reality of clinical malposition. To the best of the author's knowledge, there is no diffusion-based method that overcomes these barriers by combining: \textbf{(C1)} Non-biological rigid implant geometries (compensating for thresholding over-coverage bias, \citealp{xie2024implantsegmentation}); \textbf{(C2)} A surgically constrained placement sampler that accurately emulates the population distribution of clinical malpositions; and \textbf{(C3)} A multi-window latent diffusion synthesis module \citep{wang2025adaptiveweighting} that extends the generation band ($B_{\delta}$) beyond the implant mask to explicitly model global streaking and beam hardening. This is the gap this work aims to fill.

% #####################################################################
\section{Research Question}
% #####################################################################

What improvement does a surgical admissibility-constrained (bone density and cortical containment) 3D placement sampler, a multi-window Latent Diffusion Model (LDM) conditioned on the rigid implant geometries and an extended generation band ($B_{\delta}$), have in terms of distributional physical-coherence of the synthesized pelvic osteosynthesis implants and their local artifacts and the corresponding downstream impact on bone volumetric segmentation networks on metal artifact-affected CTs, when compared with both naive copy-paste insertion and the explicit physics-based \BASESIM{} simulator.

% #####################################################################
\section{Hypothesis}
% #####################################################################

The synthetic generation of medical anomalies and metallic implants using a rigid geometry and extended spatial boundaries, HU-conditioned diffusion model and whose placement is surgically enforced via a density-based placement sampler is more physically-coherent than unconstrained generative baseline methods and on the same footing as an explicit physical simulation (\BASESIM{}). This physical-coherence, quantified by novel Surgical Admissibility of Placement (SAP) and Boundary Feature Coherence (BFC) metrics, emulates the true clinical distribution of real-world implant placement. Subsequently, the downstream peri-implant bone volumetric segmentation (Dice, HD95) using this type of synthetic data as a training augmentation improves downstream on CTPelvic1K's test set, with full ablations revealing the non-redundant role of the multi-window encoding and extended boundary generation.

% #####################################################################
\section{Objectives}
% #####################################################################

\textbf{General Objective.} Design, implement and evaluate a physics-and-anatomy-aware latent diffusion pipline with rigid surgical constraints and extended HU-based conditioning, capable of generating physically coherent metallic implants for the improved robustness of downstream segmentation networks.

\begin{enumerate}
  \item \textbf{Validation of a multi-window representation (Compulsory Go/No-Go):} Ensure HU reconstruction error (HU to multi-window to VAE to HU again) within the bone tissue below a 25 HU threshold (Mean Absolute Error (MAE)). This will serve as a guarantee that the intensity gradients that drive downstream applications are preserved.
  \item \textbf{Design of surgically constrained placement sampler:} propose position-orientation 3D pose constrained by bone density maps \citep{arand2019pelvicring}, cortical boundaries, and corridor continuity that results in a distribution of placements that emulates the true clinical 31--60\% real-world imperfection rate, rather than a single artificial "optimized" trajectory.
  \item \textbf{Design of a conditioned synthesizer:} 2.5D ControlNet-guided Latent Diffusion Model (Stable Diffusion 1.5 backbone). A crucial detail in the formulation is that the synthesis mask must be extended with a spatial band $B_{\delta}$ ($\sim$12mm) over the implant boundary, to allow a more explicit manifestation of the streaking artifacts and beam hardening that occur outside of the implant mask in the generated image.
  \item \textbf{Formalization of physical-coherence metrics:} Formulation of SAP and BFC and Inter-slice Consistency (ISC) to decouple the evaluation of the placement sampler from the appearance synthesizer, using Wasserstein-1 distances with the real-world clinical distributions.
  \item \textbf{Evaluation of the downstream impact:} Segmentation improvement, using established metric (Dice, HD95), on real metal-impaired cases against comprehensive baselines, including the physics-based \BASESIM{}, which is verified with full per-constraint ablations.
\end{enumerate}

% #####################################################################
\section{Expected Results}
% #####################################################################

\begin{table}[h]
\centering
\small
\begin{tabularx}{\linewidth}{@{}p{3cm} p{4cm} X@{}}
\toprule
\textbf{Goal} & \textbf{Metric(s)} & \textbf{Expected Outcome} \\
\midrule
\textbf{Surgical Admissibility} \newline (Novel Sampler Eval) & \textbf{SAP}: Cortical Breach Grade \citep{smith2006iliosacral}; Density Zone fraction. & \textbf{Distribution matching:} Wasserstein-1 distance reveals that generated implant placements are close to the 31-60\% clinical real-world malposition rate, rather than unrealistic "perfect" optimization. \\
\addlinespace[4pt]
\textbf{Appearance Coherence} \newline (Synthesizer Eval) & \textbf{BFC}: Boundary Feature Coherence.\newline \textbf{ISC}: Inter-slice Consistency. \newline \textbf{Streak Amplitude}: vs real. & Lower profile discrepancy than naive copy-paste insertion, with streak amplitudes showing a statistically comparable, or better, result than the \BASESIM{} explicit physics baseline. \\
\addlinespace[4pt]
\textbf{Representation Viability} \newline (Standard Encoding Eval) & \textbf{MAE in HU} (Round-trip) & Error in the bone tissues strictly below 25 HU threshold. Log-compressed metal ranges preserve the structural topology. \\
\addlinespace[4pt]
\textbf{Downstream Detection} \newline (Clinical Utility) & \textbf{Dice, HD95} (Peri-implant restricted) & Outperforms conventional unconstrained diffusion augmentation on the audited CTPelvic1K clinical evaluation set (size pending review). \\
\addlinespace[4pt]
\textbf{Ablation validation} \newline (Constraint Verification) & All metrics above & Incremental ordering shows their non-redundant contribution. Degradation when $B_{\delta}$ band and multi-window HU encoding are removed. \\
\bottomrule
\end{tabularx}
\end{table}

% #####################################################################
\section{Datasets}
% #####################################################################

\textbf{Primary Dataset.} CTPelvic1K \citep{liu2021ctpelvic1k} will be used as the source of patient anatomy and real artifact observations. The locally available dataset6 and dataset7 volumes are inventoried separately from the full published collection; the CLINIC-metal label alone does not establish implant type or annotation availability. \emph{Strict Isolation Rule.} No patient may contribute to both training and evaluation. Test implant geometries must not enter the bank consumed by the proposed sampler or the physical simulation arm. Duplicate-volume screening and patient grouping are required to enforce this separation.

\textbf{Physical comparison protocol.} The hybrid simulation and evaluation protocol of \citet{peters2025hybrid}, also associated with the AAPM CT-MAR Challenge \citep{haneda2025aapm}, is adopted instead of an independently validated reimplementation of \BASESIM{}. Challenge data remain excluded from model training, and the scoring benchmark is reserved for final evaluation. Reproduction will document the code revision, configuration, reconstruction settings and access conditions. Both synthesis arms will use the same source anatomy and implant poses; a no-metal simulation/reconstruction control will quantify changes introduced by the physical arm alone. The published protocol is two-dimensional and evaluates MAR: extending it to pelvic osteosynthesis synthesis requires an explicit adaptation and does not inherit its validation automatically. RMSE and SSIM against the original CT will assess preservation outside the generation band, where changes are not intended; artifact realism within the affected region requires separate evaluation. No unverified numerical acceptance thresholds or MAR score calibration are adopted here.

\textbf{Local cohort audit.} Training eligibility requires manual review of three-dimensional views and multiplanar CT slices confirming absence of metal and foreign objects. HU screening only prioritizes this review, including cases from dataset6. Confirmed metal cases are candidates for testing, with pelvic osteosynthesis distinguished from other implants and external objects. Exact volume duplicates and repeated patients must be resolved before splitting; downstream evaluation additionally requires verified bone annotations. Final cohort sizes remain pending this audit.

% #####################################################################
\section{References}
% #####################################################################

\vspace{-4pt}
\scriptsize
% Lista generada con BibTeX desde refs.bib, que se regenera con
% `python scripts/build_refs.py` desde refs/clean/. No editar entradas aqui.
\renewcommand{\bibsection}{}   % el encabezado ya lo pone \section{References}
\setlength{\bibsep}{1.3pt}
\setlength{\bibhang}{8pt}
\begin{multicols}{2}
% \nocite{*} conserva las 27 entradas de refs.bib, incluidas las aun no citadas
% en el cuerpo. La lista de entradas la define la autora.
\nocite{*}
\bibliographystyle{abbrvnat}
\bibliography{../refs}
\end{multicols}
\end{document}
